% Generated by regenerate_supplementary.py from results/dataset/tables/class_stats.csv
\begin{table}[H]
\centering
\small
\setlength{\tabcolsep}{4pt}
\begin{adjustbox}{max width=\linewidth}
\begin{tabular}{lrrrrrr}
\toprule
Group & Patients & Height (cm) & Weight (kg) & Slices & In-plane (mm) & Slice step (mm) \\
\midrule
DCM & 20 & 171.1 & 75.5 & 9.8 & 1.51 & 10.0 \\
HCM & 20 & 169.8 & 80.3 & 9.4 & 1.49 & 9.8 \\
MINF & 20 & 172.2 & 74.2 & 8.7 & 1.49 & 9.8 \\
NOR & 20 & 172.2 & 77.9 & 8.5 & 1.52 & 9.8 \\
ARV & 20 & 168.9 & 66.9 & 11.2 & 1.55 & 7.3 \\
Test (n=50) & 50 & 170.0 & 80.7 & 10.8 & 1.52 & 9.1 \\
\bottomrule
\end{tabular}
\end{adjustbox}
\caption{Patients per class (means). The training set has 20 patients per class; the test set has the same number of patients per class in ACDC, but its labels are not available. Slice counts range from 6 to 21; slice steps take the values 5, 6.5, 7, 10 mm.}
\label{tab:supp-classes}
\end{table}
